% --- Back cover ---
% Large PERSONIX wordmark, no logo image. Cream paper inherited globally.
% Absolute (tikz overlay) positioning so the footer sticks to the bottom.
\clearpage
\thispagestyle{empty}
\begin{tikzpicture}[remember picture, overlay]
  % --- Wordmark ---
  \node[anchor=north, font=\sffamily\bfseries\fontsize{48}{52}\selectfont]
    at ([yshift=-26mm]current page.north) {PERSONIX};
  \node[anchor=north, font=\rmfamily\itshape\large]
    at ([yshift=-44mm]current page.north) {Cambiamento senza compromessi};
  \node[anchor=north, font=\rmfamily\small,
        text width=\paperwidth-50mm, align=center]
    at ([yshift=-54mm]current page.north)
    {Una rete di reputazione decentralizzata, incensurabile e incorruttibile};
  % --- Body ---
  \node[anchor=north, text width=\paperwidth-50mm, align=justify, font=\small]
    at ([yshift=-72mm]current page.north)
    {%
      Lo Stato funzionava quando la comunicazione era lenta, le decisioni erano
      locali e l'autorità viveva nella stessa città che governava. Le reti di
      oggi ribaltano tutti e tre questi presupposti --- e le istituzioni
      costruite su di essi non si adattano più al mondo che governano. Questo
      libro descrive uno strumento che invece vi si adatta: una rete di
      reputazione decentralizzata in cui l'identità è tua e resta tua, la
      verità è verificabile pubblicamente e la corruzione è un suicidio
      reputazionale.

      \smallskip
      Non un manifesto. Non una previsione. Un progetto operativo su come si
      può costruire il successore dello Stato --- volontariamente, una
      dichiarazione alla volta.
    };
  % --- Footer (anchored to the bottom of the page) ---
  \node[anchor=south, font=\sffamily\footnotesize,
        text width=\paperwidth-50mm, align=center]
    at ([yshift=12mm]current page.south)
    {%
      Pavel Kudrna\quad\textbullet\quad Praga, 2026\par
      \href{https://personix.org}{personix.org}\par
      Distribuito sotto licenza Creative Commons Attribution-ShareAlike 4.0
      International (CC BY-SA 4.0).
    };
\end{tikzpicture}%
\mbox{}%
\clearpage
